
\section{Computation of the mass operator}

\SourcePage{576}{581}

Using the example of the computation of the mass operator, we shall demonstrate the method of direct regularisation of Feynman integrals.

In the first non-vanishing approximation, the mass operator is represented by the loop in the diagram

% [Feynman diagram (119.1): first-order mass operator — electron line carrying momentum p enters from the left, emits a virtual photon line carrying momentum k (forming a loop with internal electron line carrying momentum p-k), and exits to the right with momentum p]

\begin{equation}
\label{eq:119.1}
\end{equation}

\noindent The corresponding integral is
\[
-i\bar{\mathcal{M}}(p) = (-ie)^2 \int \gamma^\mu G(p-k)\gamma^\nu D_{\mu\nu}(k)\,\frac{d^4k}{(2\pi)^4};
\]
substituting the propagators and bringing together the factors $\gamma^\mu \dots \gamma_\mu$ with the help of the formulas (22.6), we obtain
\begin{equation}
\bar{\mathcal{M}}(p) = -\frac{8\pi i}{(2\pi)^4}\,e^2 \int \frac{2m - (\gamma p) + (\gamma k)}{[(p-k)^2 - m^2](k^2 - \lambda^2)}\,d^4k
\label{eq:119.2}
\end{equation}
(by the bar over $\mathcal{M}$ we denote the unregularised value of the integral). A fictitious ``photon mass'' $\lambda$ has been introduced into the photon propagator in order to eliminate the infrared divergence (as in \S\,117).

We transform the integral with the help of formula (131.4), taking the two factors in the denominator of~(\ref{eq:119.2}) as $a_1$ and $a_2$. After a simple regrouping of the terms in the denominator of the new integral, we obtain
\begin{equation}
\bar{\mathcal{M}}(p) = -\frac{8\pi i}{(2\pi)^4}\,e^2 \int d^4k \int_0^1 dx\,\frac{2m - (\gamma p) + (\gamma k)}{[(k - px)^2 - a^2]^2},
\label{eq:119.3}
\end{equation}
where
\begin{equation}
a^2 = m^2 x^2 - (p^2 - m^2)x(1-x) + \lambda^2(1-x).
\label{eq:119.4}
\end{equation}

The substitution $k \to k + px$ brings the integrand in~(\ref{eq:119.3}) to a form in which the denominator depends only on $k^2$. In doing so, however, according to (131.17), (131.18), an additive constant is added to the integral:
\begin{equation}
\bar{\mathcal{M}}(p) = -\frac{8\pi i}{(2\pi)^4}\,e^2 \bigg\{\int d^4k \int_0^1 dx\,\frac{2m(1-x) - \frac{i\pi^2}{4}(\gamma p)}{(k^2 - a^2)^2} \bigg\},
\label{eq:119.5}
\end{equation}
the term with $(\gamma k)$ in the numerator being now omitted as it vanishes upon integration over the directions of the 4-vector $k$ (cf.\ (131.8)).

\SourcePage{577}{582}

Regularisation of this integral consists in such subtractions as would reduce it to an expression of the form (110.20). The latter vanishes upon multiplication by the wave amplitude $u(p)$, if $p$ is the 4-momentum of a real electron. Without introducing $u(p)$ explicitly, one can formulate this condition as the requirement that $\mathcal{M}(p)$ vanish upon the substitution
\begin{equation}
(\gamma p) \to m, \qquad p^2 \to m^2.
\label{eq:119.6}
\end{equation}
The form of the integral~(\ref{eq:119.5}) is convenient in that the 4-vector $p$ enters it only in the form $\gamma p$ and $p^2$ (and terms of the form $kp$ are absent).

Subtracting from~(\ref{eq:119.5}) the same expression with the substitution~(\ref{eq:119.6}), we obtain
\begin{multline}
-\frac{8\pi i e^2}{(2\pi)^4} \bigg\{ \int d^4k \int_0^1 dx \cdot [2m - (\gamma p)(1-x)] \bigg[\frac{1}{(k^2 - a^2)^2} - \frac{1}{(k^2 - a_0^2)^2}\bigg] \\
- \int d^4k \int_0^1 dx\,\frac{1-x}{(k^2 - a_0^2)^2}\,(\gamma p - m) - \frac{i\pi^2}{4}\,(\gamma p - m) \bigg\},
\label{eq:119.7}
\end{multline}
where
\[
a_0^2 = m^2 x^2 + \lambda^2(1-x).
\]

For the final regularisation, however, one more subtraction must be performed: according to (110.20), upon the substitution~(\ref{eq:119.6}) not only $\mathcal{M}(p)$ as a whole should vanish, but so should the same expression without one of the factors $\gamma p - m$. The corresponding subtractions eliminate entirely the second and third terms in the curly brackets of~(\ref{eq:119.7})\footnote{In doing so, during the process of ``renormalisation on the fly'' (cf.\ p.\,545), we drop corrections to the renormalisation constant $Z_1$ (see \S\,110). The corresponding integrals diverge logarithmically. If one introduces a ``cut-off parameter'' $\Lambda^2 \gg m^2$, $p^2$, restricting the region of integration over $d^4k$ by the condition $k^2 \leqslant \Lambda^2$, this correction can be computed explicitly. The computation gives the result
\begin{equation}
Z_1 = 1 + Z_1^{(1)}, \quad Z_1^{(1)} = -\frac{\alpha}{2\pi}\bigg[\frac{1}{2}\ln\frac{\Lambda^2}{m^2} + \ln\frac{\lambda^2}{m^2} + \frac{9}{4}\bigg].
\label{eq:119.7a}
\end{equation}}.
The first integral is preliminarily transformed by introducing one more auxiliary integration with the help of formula (131.5), putting in it $n=2$ and taking $k^2 - a^2$ and $k^2 - a_0^2$ as $a$ and $b$. Then expression~(\ref{eq:119.7}) takes the form

\SourcePage{578}{583}

\[
(\gamma p - m)\frac{16\pi i}{(2\pi)^4}\,e^2 \int d^4k \int_0^1 dx \int_0^1 dz\,\frac{[2m - (\gamma p)(1-x)]x(1-x)}{[k^2 - a_0^2 + (p^2 - m^2)x(1-x)z]^3}
\]
(here we have also used the identity $p^2 - m^2 = (\gamma p - m)(\gamma p + m)$). We immediately perform the integration over $d^4k$. Assuming that $p^2 - m^2 < 0$, and using (131.14), we obtain
\[
(\gamma p - m)\frac{e^2}{2\pi} \int_0^1 dx \int_0^1 dz\,\frac{(\gamma p + m)[2m - (\gamma p)(1-x)]x(1-x)}{m^2 x^2 + \lambda^2(1-x) + (m^2 - p^2)x(1-x)z}.
\]

It now remains to subtract the factor $(\gamma p - m)$, temporarily omitted, and to perform the same integral with the substitution~(\ref{eq:119.6}); after straightforward manipulations we obtain
\begin{multline}
\mathcal{M}(p) = (\gamma p - m)\frac{e^2}{2\pi} \times \\
\times \int_0^1 dx \int_0^1 dz\,\frac{m(1-x^2) - (\gamma p)(1-x)^2}{m^2 x + (m^2 - p^2)(1-x)z}\bigg[1 - \frac{2x(1+x)}{x^2 + (\lambda/m)^2}\bigg]
\label{eq:119.8}
\end{multline}
(in the common denominator the term with $\lambda^2$ is omitted, since this does not lead to divergences here; in another place $\lambda^2(1-x)$ is replaced by

\SourcePage{579}{584}

\noindent $\lambda^2$, since the infrared divergence will correspond to the divergence at $x \to 0$).

Integration in~(\ref{eq:119.8}) (first over $z$, then over $x$) is rather lengthy, but elementary and leads to the following final result:
\begin{multline}
\mathcal{M}(p) = \frac{\alpha}{2\pi m}\,(\gamma p - m)^2 \bigg\{\frac{1}{2(1-\rho)}\bigg(1 - \frac{2-3\rho}{1-\rho}\ln\rho\bigg) - \\
- \frac{\gamma p + m}{m\rho}\bigg[\frac{1}{2(1-\rho)}\bigg(2 - \rho + \frac{\rho^2 + 4\rho - 4}{1-\rho}\ln\rho\bigg) + 1 + 2\ln\frac{\lambda}{m}\bigg]\bigg\},
\label{eq:119.9}
\end{multline}
where we have denoted
\[
\rho = \frac{m^2 - p^2}{m^2}.
\]
(\emph{R.~Karplus, N.\,M.~Kroll}, 1950). The integral is computed under the assumption $\rho > 0$, and moreover $\rho \gg \lambda/m$. In accordance with the rule for going around poles, the analytic continuation of expression~(\ref{eq:119.9}) into the region $\rho < 0$ is defined by the substitution $m \to m - i0$; in doing so $\rho \to \rho - i0$, so that $\ln\rho$ for $\rho < 0$ must be understood as
\begin{equation}
\ln\rho = \ln|\rho| - i\pi, \quad \rho < 0.
\label{eq:119.10}
\end{equation}

Let us examine the behaviour of the mass operator for $p^2 \gg m^2$. We then have $-\rho \approx p^2/m^2 \gg 1$, and with logarithmic accuracy
\begin{equation}
\mathcal{M}(p) = -[\mathcal{G}^{-1}(p) - G^{-1}(p)] \approx \frac{\alpha}{4\pi}\,(\gamma p)\ln\frac{p^2}{m^2}.
\label{eq:119.11}
\end{equation}
As in the case of the photon propagator (cf.\ formulas (113.15), (113.16) for the polarisation operator), the correction to $G^{-1}$ turns out to be small only at not too large energies, namely when
\[
\frac{\alpha}{4\pi}\ln\frac{p^2}{m^2} \ll 1.
\]

In the present case, however, the logarithmic growth is in a certain sense fictitious: it can be eliminated by an appropriate choice of gauge, i.e.\ of the function $D^{(l)}$ in the photon propagator (\emph{L.\,D.~Landau, A.\,A.~Abrikosov, I.\,M.~Khalatnikov}, 1954). Namely, one must set (in the notation of \S\,103)
\begin{equation}
D^{(l)} = 0,
\label{eq:119.12}
\end{equation}
whereas formula~(\ref{eq:119.9}) was obtained in the gauge
\begin{equation}
D^{(l)} = D.
\label{eq:119.13}
\end{equation}
This property of the gauge~(\ref{eq:119.12}) makes it especially convenient for investigating the behaviour of the theory at $p^2 \gg m^2$, which will be used below, in \S\,132.

\SourcePage{580}{585}

To prove the assertion made, we note that if we are interested only in terms $\sim e^2$, then the transformation from the gauge~(\ref{eq:119.13}) to the gauge~(\ref{eq:119.12}) can be regarded as infinitesimally small. Correspondingly, one can directly use formula (105.14), putting in it
\[
d^{(l)}(q) = -\frac{D}{q^2} = -\frac{4\pi}{(q^2)^2},
\]
and also replacing, with the required accuracy, the function $\mathcal{G}$ in the integrand by $G$. In the integral over $d^4q$, the essential region is $q \gg p$; in this case $G(p-q)$ in the integrand is much smaller than $G(p)$, and it can be neglected. Then
\[
\delta\mathcal{G}^{-1} = -\mathcal{G}^{-2}(p)\,\delta\mathcal{G}(p) = -ie^2\,\mathcal{G}^{-1}(p) \int d^{(l)}(q)\,\frac{d^4q}{(2\pi)^4}.
\]

Finally, applying the transformation (113.11), (113.12), we obtain
\[
\delta\mathcal{G}^{-1} = -\frac{e^2}{4\pi}\,\mathcal{G}^{-1}(p) \int \frac{d(-q^2)}{-q^2} \approx -\frac{e^2}{4\pi}\,(\gamma p)\ln\frac{\Lambda^2}{p^2},
\]
where $\Lambda$ is an auxiliary upper limit, the divergence at which is removed by renormalisation. The latter consists in subtracting the same expression evaluated at $p^2 \approx m^2$, so that we finally have
\[
\delta\mathcal{G}^{-1} = \frac{e^2}{4\pi}\,(\gamma p)\ln\frac{p^2}{m^2}.
\]
This expression exactly cancels the difference $\mathcal{G}^{-1} - G^{-1}$ from~(\ref{eq:119.11}).

Finally, let us address the question of the reasons necessitating the introduction of a finite ``photon mass'' $\lambda$ for the regularisation of the integral~(\ref{eq:119.2}), closely connected with its behaviour at $p^2 \to m^2$.

First of all, let us note that this integral is by itself finite at $\lambda = 0$ (for the elimination of the inessential divergence in this respect at large $k$, we assume that the integral is taken over a large but finite region of $k$-space). The need for introducing $\lambda$ arises from subtracting the renormalisation integral, which would otherwise diverge at $p^2 = m^2$. Let us therefore clarify how the unregularised mass operator behaves as $p^2 \to m^2$. Since this behaviour depends significantly on the choice of gauge (whereas the integral~(\ref{eq:119.2}) was written already for a specific choice---(\ref{eq:119.13})), let us consider the general case of an arbitrary gauge.

We again make use of the transformation (105.14). Representing

\SourcePage{581}{586}

\noindent $d^{(l)}$ in the form
\begin{equation}
d^{(l)}(q) = -\frac{\delta D^{(l)}}{q^2} = -\frac{4\pi}{(q^2)^2}\,\delta a(q^2),
\label{eq:119.14}
\end{equation}
we shall consider $\delta a$ to be a variation of the function $a(q^2)$, substantially changing only over intervals $q^2 \sim m^2$ and finite at $q^2 \approx m^2$. In the integrand on the right side of (105.14), in the difference $\mathcal{G}(p) - \mathcal{G}(p-q)$, for small $q$ both terms are close and the integral converges. Since for small $q$
\[
\mathcal{G}(p-q) \sim \frac{1}{p^2 - m^2 - 2pq},
\]
$\mathcal{G}(p-q)$ can be omitted in comparison with $\mathcal{G}(p)$ when $q \gg (p^2 - m^2)/m$. The integral
\[
\delta\mathcal{G}(p) = ie^2\,\mathcal{G}(p) \int d^{(l)}(q)\,\frac{d^4q}{(2\pi)^4} = -\frac{e^2}{4\pi}\,\mathcal{G}(p) \int \delta a(q^2)\,\frac{d(-q^2)}{-q^2}
\]
diverges logarithmically in the region
\[
(p^2 - m^2)/m^2 \ll q^2 \ll m^2.
\]
With logarithmic accuracy we therefore have
\[
\frac{\delta\mathcal{G}}{\mathcal{G}} = -\frac{e^2}{4\pi}\,\delta a(m^2)\ln\frac{m^2}{p^2 - m^2}.
\]

This equation can be integrated. Noting that when $\alpha^2 \equiv e^2 \to 0$ the exact propagator $\mathcal{G}$ must coincide with the propagator of free particles $G$, we obtain
\begin{equation}
\mathcal{G}(p) = \frac{1}{\gamma p - m}\bigg(\frac{m^2}{p^2 - m^2}\bigg)^{\frac{\alpha}{2\pi}(C - a_0)},
\label{eq:119.15}
\end{equation}
where $a_0 = a(m^2)$, $C$ is a certain constant. To determine the latter, we compare the expression
\begin{equation}
\mathcal{G}^{-1}(p) = (\gamma p - m)\bigg[1 + \frac{\alpha}{2\pi}(C - a_0)\ln\rho\bigg],
\label{eq:119.16}
\end{equation}
obtained from~(\ref{eq:119.15}) in the first approximation with respect to $\alpha$, with the analogous expression obtained from the integral~(\ref{eq:119.2}) at $\lambda = 0$\footnote{To obtain~(\ref{eq:119.16}), there is no need to redo the computation from scratch. The term $\sim \ln\rho$ in~(\ref{eq:119.9}) was in fact obtained under the assumption $\rho \gg \lambda$, which permits the passage to $\lambda \to 0$. The term $\sim \ln(\lambda/m)$ arises from the subtraction of the renormalisation integral and is absent in the original integral~(\ref{eq:119.2}).}:
\begin{equation}
\mathcal{G}^{-1}(p) = (\gamma p - m)\bigg[1 + \frac{\alpha}{\pi}\ln\rho\bigg].
\label{eq:119.17}
\end{equation}
According to the definition~(\ref{eq:119.14}), the function $a(q^2)$ coincides with the ratio $D^{(l)}/D$. Therefore the gauge~(\ref{eq:119.13}), to which~(\ref{eq:119.17}) pertains, corresponds to $a = a_0 = 1$. Requiring the coincidence of~(\ref{eq:119.16}) and~(\ref{eq:119.17}) at this value of $a_0$, we find $C = 3$.

\SourcePage{582}{587}

Thus we finally find the following limiting expression (\emph{infrared asymptotic}) of the unrenormalised electron propagator as $p^2 \to m^2$:
\begin{equation}
\mathcal{G}(p) = \frac{\gamma p + m}{p^2 - m^2}\bigg(\frac{m^2}{p^2 - m^2}\bigg)^{\frac{\alpha}{2\pi}(3 - a_0)}.
\label{eq:119.18}
\end{equation}
(\emph{A.\,A.~Abrikosov}, 1955). Let us stress that the validity of this formula is subject to the inequalities $\alpha \ll 1$, $|\ln\rho| \gg 1$, whereas the formulas of perturbation theory would require the condition $\alpha|\ln\rho|/2\pi \ll 1$ as well. Let us also note that the sign of the difference $p^2 - m^2$ is immaterial here, since the imaginary part of expression~(\ref{eq:119.18}) would in any case be beyond its limits of accuracy.

The renormalised propagator must have a simple pole at $p^2 = m^2$. We see that~(\ref{eq:119.18}) satisfies this requirement only in the gauge in which
\begin{equation}
D^{(l)} = 3D.
\label{eq:119.19}
\end{equation}
(so that $a_0 = 3$). In this case the regularisation of the Feynman integral (having the purpose of eliminating its divergence at the upper limits) will not require the introduction of a finite ``photon mass''. In other gauges, however, a vanishing photon mass leads to the appearance of a branch point at $p^2 = m^2$ instead of a simple pole, and the removal of this ``defect'' requires the introduction of a finite parameter $\lambda$.
